\documentclass{article}

%% Page Margins %%
\usepackage{geometry}
\usepackage{colortbl}   % for \rowcolor
\usepackage{xcolor}     % for \textcolor and [HTML] colour spec
\usepackage{array}      % for p{} column type
\geometry{
    top = 0.75in,
    bottom = 0.75in,
    right = 0.75in,
    left = 0.75in,
}

\usepackage{amsmath}
\usepackage{graphicx}
\usepackage{parskip}

\title{Assembly Project: COLUMNS}

% TODO: Enter your name
\author{Faiyad Ahmed Masnoon}

\begin{document}
\maketitle

\section{Instruction and Summary}

\begin{enumerate}

    \item Which milestones were implemented? \\
    % TODO: List the milestone(s) and in the case of 
    %       Milestones 4 & 5, list what features you 
    %       implemented, sorted into easy and hard 
    %       categories.
    \textbf{Answer: } Milestones 1, 2, 3, 4 and 5 (6 easy features + 1 hard). 

    \item How to view the game:
    % TODO: specify the pixes/unit, width and height of 
    %       your game, etc.  NOTE: list these details in
    %       the header of your breakout.asm file too!
    
    \begin{enumerate}

    \item Set bitmap size to 256x256
    \item Set units to 8 for both height and width
    \item Set base address for display: 0x10008000 (\$gp)
    \item Run the columns.asm file (I used Saturn IDE).

    \end{enumerate}

\item Game Summary:
% TODO: Tell us a little about your game.
\begin{itemize}
\item Implementation of the Columns game in assembly!
\item The game generates a 3x1 column at the top of the screen, that falls down. If there is a match, the gems disappear. If any gem touches the top of the playable area, the game ends.\\Also has score counter, gravity, ghost outline, difficulty select and more!
\end{itemize}

    
\end{enumerate}

% \section{Attribution Table}
% % TODO: If you worked in partners, tell us who was 
% %       responsible for which features. Some reweighting 
% %       might be possible in cases where one group member
% %       deserves extra credit for the work they put in.

% \begin{center}
% \begin{tabular}{|| c | c ||}
% \hline
%  Student 1 (Name and student number) &  Student 2 (Name and student number) \\ 
%  \hline
%  Task & Task\\
%  \hline
%  Task & Task\\
%  \hline
%  Task & Task\\ 
%  \hline
%  Task & Task\\ 
%  \hline
%  Task & Task\\
%  \hline
%  Task & Task\\  
%  \hline
% \end{tabular}
% \end{center}

% TODO: Fill out the remainder of the document as you see 
%       fit, including as much detail as you think 
%       necessary to better understand your code. 
%       You can add extra sections and subsections to 
%       help us understand why you deserve marks for 
%       features that were more challenging than they
%       might initially seem.

\section{How to Play}
\begin{enumerate}
    \item Input 1, 2 or 3 depending on how difficult you want the game to be. 1: easy; 2: medium; 3:hard. Do note that the difficulty increases over time regardless of initial choice. Choosing medium or hard difficulty will provide a score multiplier.
    \item A column will be spawned near the top. Use A/D to move the column left or right. 
    \item Use S to move the column downwards.
    \item Use W to shuffle the order of the gems in the column.
    \item There will be an outline to show where the column will land.
    \item Press P to pause the game.
    \item If you run out of space, the game will end. Press R to retry or Q to exit the game completely.
\end{enumerate}

\newpage

\section{Extra Details}
    \begin{figure}[ht!]
        \centering
        \includegraphics[width=0.3\textwidth]{state.png}
        \caption{Columns, showing the score counter, pause symbol and ghost outline!}
        \label{Instructions}
    \end{figure}

    \textbf{Easy features implemented (6):}
    \begin{enumerate}
        \item \textbf{Gravity (\#1): } Spawned column will move down by itself.
        \item \textbf{Auto increase in gravity (\#2): } Over time, the column will fall faster, increasing the difficulty. This is set to increase every 10s (every 600 frames). \\
        To change this you can edit line 182.
        \item \textbf{Difficulty selection (\#3): } Once you start the game, you can choose the difficulty. Enter 1, 2 or 3.
        \begin{figure}[ht!]
            \centering
            \includegraphics[width=0.2\textwidth]{difficulty.png}
            \caption{Difficulty select screen}
            \label{Instructions}
        \end{figure}
        \item \textbf{Game Over + Retry (\#4): } Once the new column has no place to spawn (already occupied by previous gems), the game ends. The screen turns red and there is a big R in the screen, representing the option to Retry. You can press R to retry, or Q to quit.\\
        Retrying resets the score, and starts a new game.
        \begin{figure}[ht!]
            \centering
            \includegraphics[width=0.2\textwidth]{retry.png}
            \caption{Game over screen}
            \label{Instructions}
        \end{figure}
        \item \textbf{Pause (\#6): } Press P to pause the game. A pause sign appears to indicate the game is paused. Press P again to resume.
        \item \textbf{Ghost Outline (\#7): } A gray outline shows where the current column will fall.
    \end{enumerate}
    \textbf{Hard features implementd (1): }
    \begin{enumerate}
        \item \textbf{Score (\#1): } A score counter that increases whenever you make matches. Harder difficulty has a score multiplier. Also when you have chain matches, there is a multiplier. Max score is 999.
    \end{enumerate}


\section{Memory Layout}
\begin{table}[h]
\centering
\begin{tabular}{|l|l|l|p{6cm}|}
\hline
\rowcolor[HTML]{2E4057}
\textcolor{white}{\textbf{Label}} &
\textcolor{white}{\textbf{Type / Size}} &
\textcolor{white}{\textbf{Initial Value}} &
\textcolor{white}{\textbf{Purpose}} \\
\hline
\texttt{ADDR\_DSPL}    & \texttt{.word (4 B)}    & \texttt{0x10008000}  & Base address of bitmap display MMIO \\
\texttt{ADDR\_KBRD}    & \texttt{.word (4 B)}    & \texttt{0xFFFF0000}  & Base address of keyboard MMIO \\
\texttt{KEY\_A/D/../2/3}& \texttt{.word ×10}       & \texttt{0x61...0x33} & ASCII codes for each control key \\
\texttt{FIELD\_X\_MIN} & \texttt{.word (4 B)}    & \texttt{11}          & Leftmost valid field column (x) \\
\texttt{FIELD\_X\_MAX} & \texttt{.word (4 B)}    & \texttt{20}          & Rightmost valid field column (x) \\
\texttt{FIELD\_Y\_MAX} & \texttt{.word (4 B)}    & \texttt{20}          & Bottom-most valid field row (y) \\
\texttt{SLEEP\_MS}     & \texttt{.word (4 B)}    & \texttt{16}          & Frame sleep duration\\
\texttt{FALL\_DELAY}   & \texttt{.word (4 B)}    & \texttt{40}          & Frames between automatic drops\\
\texttt{FALL\_DELAY\_MID}   & \texttt{.word (4 B)}    & \texttt{25}          & FALL\_DELAY for medium diff\\
\texttt{FALL\_DELAY\_HRD}   & \texttt{.word (4 B)}    & \texttt{10}          & FALL\_DELAY for hard diff\\
\texttt{COL\_BG}       & \texttt{.word (4 B)}    & \texttt{0x111111}    & Background colour (dark grey) \\
\texttt{COL\_WALL}     & \texttt{.word (4 B)}    & \texttt{0x888888}    & Wall and floor colour (grey) \\
\texttt{COL\_FIELD}    & \texttt{.word (4 B)}    & \texttt{0x222222}    & Empty field cell colour \\
\texttt{gem\_colors}   & \texttt{.word ×6 (24 B)}& \texttt{0xFF0000...} & RGB values for gem colours indexed 0--5 \\
\texttt{SCORE}   & \texttt{.word (4 B)}& \texttt{0} & Initial Score \\
\texttt{SCORE\_MP}   & \texttt{.word (4 B)}& \texttt{1} & Score multiplier for easy difficulty \\
\texttt{SCORE\_MP\_MID}   & \texttt{.word (4 B)}& \texttt{2} & Score multiplier for medium difficulty \\
\texttt{SCORE\_MP\_HRD}   & \texttt{.word (4 B)}& \texttt{3} & Score multiplier for hard difficulty \\
\texttt{CHAIN}   & \texttt{.word (4 B)}& \texttt{1} & Chain multiplier\\
\texttt{TIMER}   & \texttt{.word (4 B)}& \texttt{0} & Initial time, for auto\_diff \\

\hline
\end{tabular}
\caption{Memory used}
\label{tab:data-layout}
\end{table}

    



\end{document}